\section{Discussion and limitations}\label{sec:disc}
\paragraph{What is score-free and what is not.} The finite-sample marginal coverage law is score-free (Proposition~\ref{prop:beta}, Fig.~\ref{fig:beta}); it is the right first test of any evaluation pipeline, and it fails to discriminate between scores. Everything a user may care about beyond it, namely region size, coverage across the simplex and robustness to shift, depends on the score geometry, the noise law and the calibration scheme, and on classifier outputs several of the simplex scores collapse to the same sets (Proposition~\ref{prop:vertex}, Table~\ref{tab:e2}). Conditional-coverage metrics computed from bins or slabs have a null spread (Fig.~\ref{fig:null}) that must be reported next to any comparison.

\paragraph{Scope of the conclusions.} Results hold for exchangeable data; for the tested dimensions $K\in\{3,5,10\}$ and calibration sizes; for the generators described in Section~\ref{sec:exp} (fixed regression function per configuration, Gaussian covariates, two predictors); and for one CNN on CIFAR-10 and its CIFAR-10H annotations. They should not be extrapolated to (i) $K>10$ or to Lebesgue-volume comparisons for $K\ge 10$, where we use a different, dimension-valid measure (mass of a reference law) and compare scores only within it; (ii) full conformal or jackknife+ with neural predictors, which we did not run (CV+ is run only for the Euclidean score with the linear predictor); (iii) covariate or label shift beyond the single family in Section~\ref{sec:exp}; (iv) batch multivalid or other group-conditional methods \citep{jung2023batch}, which we did not run; (v) response distributions with heavy mass on the boundary other than those generated, in particular the zero-heavy CIFAR-10H proportions, where ties give only a lower bound.

\paragraph{Limitations.} We claim no new coverage guarantee and no conditional guarantee: exact conditional coverage is unattainable without assumptions, and the oracle measurements are specific to the generators. The delta-method prediction is first-order and fails for the Aitchison score at the tested precisions (Fig.~\ref{fig:delta}). The Aitchison and Dirichlet-likelihood scores depend on the zero-replacement constant $\varepsilon$, which we varied only in a sensitivity run. Efficiency depends on the reference measure (Lebesgue on the simplex, or the marginal law of $y$); rankings under the two may differ and we do not compare across them. RAPS is not tuned. Intervals reflect variability over data and calibration draws for fixed regression functions; they do not account for variation over the choice of regression function, since each configuration fixes one with seed $0$. The synthetic worlds were designed by us; their noise levels (e.g.\ $\phi_0=20$) determine how sharp the differences between scores are. Our novelty claim is limited to the literature we read (Section~\ref{sec:related}); we did not search exhaustively.
